
        \documentclass[10pt,brazil,a4paper]{exam}
        \usepackage[latin1]{inputenc}
        \usepackage[portuguese]{babel}
        \usepackage[dvips]{graphicx}
        %\usepackage{multicol}
        %\usepackage{shadow}
        %\usepackage{pifont}
        %\usepackage{listings}
        %\usepackage{fancyvrb}
        
        \newcommand*\varhrulefill[1][0.4pt]{\leavevmode\leaders\hrule height#1\hfill\kern0pt}
        
        \def\drawLines#1{{\color{cyan}\foreach \x in {1,...,#1}{\par\vspace{2mm}\noindent\hrulefill}}}
        
        \usepackage{enumitem}
        \usepackage{multirow}
        \usepackage{amsmath}
        \usepackage{changepage,ifthen}
        %\usepackage{boxedminipage}
        %\usepackage{theorem}
        \usepackage{verbatim}
        \usepackage{tabularx}
        %\usepackage{moreverb}
        \usepackage{times}
        %\usepackage{relsize}
        \usepackage{pst-barcode}
        \usepackage{tikz}
        \setlength{\textwidth}{185mm}
        \setlength{\oddsidemargin}{-0.5in}
        \setlength{\evensidemargin}{0in}
        \setlength{\columnsep}{8mm}
        \setlength{\topmargin}{-28mm}
        \setlength{\textheight}{265mm}
        \setlength{\itemsep}{0in}
        \begin{document}
        \pagestyle{empty}
        %\lstset{language=python}
        \makeatletter\renewcommand*\cleardoublepage{\ifodd\c@page \else\hbox{}\newpage\fi}
\makeatother
\cleardoublepage
\begin{table}[h]\centering
\begin{tabular}{|p{16mm}|p{16cm}|}
\hline\multirow{4}{*}{\hspace{-2mm}\includegraphics[width=2cm]{../../../figs/ufabc.eps}}
&\vspace{-2mm}\noindent\large\textbf{Universidade Federal do ABC}\\
&\noindent\textbf{Course:} Processamento da Informa��o - BC0505\\
&\noindent\textbf{Teacher(s):} Francisco de Assis Zampirolli\\
&\noindent\textbf{Period:} 2/2016\hfill\textbf{Modality:} Presencial\hfill\textbf{Date:} 22/06/2016\hspace{-8mm}\\
\hline\end{tabular}
\end{table}
\vspace{-4mm}\small{
\noindent\textbf{Student:}  Fulano Junior
\hfill\textbf{Registration:} 11000123
\hfill\textbf{Room:} 2015$\_$BC0505$\_$q3$\_$A2
\hspace{-1mm}

\vspace{8mm}
\noindent\textbf{Sig.:}\rule{11.5cm}{0.1pt}\hfill\hspace{1cm}
}\begin{pspicture}(6,0in)
\psbarcode[scalex=1.6,scaley=0.35]{000011000123}{}{ean13}
\end{pspicture}
\\{\scriptsize

\noindent\textbf{Instructions: }\vspace{-1mm}\begin{verbatim}

1. Paint only INSIDE OF THE CIRCLES of each question.  2. No Rasure. 3. Each question has only one correct answer.
4. It will be considered only the answers in "Part 1" area on this page to the questions of multiple choice.
\end{verbatim}
\begin{center}\textbf{Parte 1 - Answer Sheet - Do not use this sheet as a draft!}\end{center}
\vspace{-5mm}\noindent\varhrulefill[0.4mm]\vspace{1mm}

\vspace{-3mm}\noindent\varhrulefill[0.4mm]\vspace{1mm}

\begin{center}
\begin{tikzpicture}[font=\tiny]
  \foreach \letter/\position in {A/1,B/2,C/3,D/4,E/5} {
    \node[inner sep=3pt] at ({\position * 0.5},0) {\letter};
  }
  \foreach \line in {1,...,9} {
     \begin{scope}[xshift=0cm,yshift=-(\line-1+1)*5mm]
       \foreach \letter/\position in {A/1,B/2,C/3,D/4,E/5} {
           \node at (-0.1,0) {\line};
           \node[fill=black!30,draw,circle,inner sep=3pt] at ({\position * 0.5},0) {};
           \node[fill=white,draw,circle,inner sep=2pt] at ({\position * 0.5},0) {};
       }
     \end{scope}
  }
\end{tikzpicture}\hspace{1 cm}


\end{center}
\vspace{1cm}\noindent\varhrulefill[0.4mm]\vspace{1mm}

\vspace{-3mm}\noindent\varhrulefill[0.4mm]\vspace{1mm}


\begin{table}[h]
\centering
\textbf{Calculation of concepts} \\ \vspace{5mm}
\begin{tabular}{|c|c|c|c|} \hline
Questions 1-12 & Textual Questions & Final\\ \hline
& &  \\
& &  \\
& &  \\ \hline
\end{tabular}
\end{table}\newpage\thispagestyle{empty}\mbox{}
 \ \ 
\newpage
\makeatletter\renewcommand*\cleardoublepage{\ifodd\c@page \else\hbox{}\newpage\fi}
\makeatother
\cleardoublepage
\begin{table}[h]\centering
\begin{tabular}{|p{16mm}|p{16cm}|}
\hline\multirow{4}{*}{\hspace{-2mm}\includegraphics[width=2cm]{../../../figs/ufabc.eps}}
&\vspace{-2mm}\noindent\large\textbf{Universidade Federal do ABC}\\
&\noindent\textbf{Course:} Processamento da Informa��o - BC0505\\
&\noindent\textbf{Teacher(s):} Francisco de Assis Zampirolli\\
&\noindent\textbf{Period:} 2/2016\hfill\textbf{Modality:} Presencial\hfill\textbf{Date:} 22/06/2016\hspace{-8mm}\\
\hline\end{tabular}
\end{table}
\vspace{-4mm}\small{
\noindent\textbf{Student:}  Fulano Junior
\hfill\textbf{Registration:} 11000123
\hfill\textbf{Room:} 2015$\_$BC0505$\_$q3$\_$A2
\hspace{-1mm}

\vspace{8mm}
\noindent\textbf{Sig.:}\rule{11.5cm}{0.1pt}\hfill\hspace{1cm}
}
\vspace{4mm}
\\{\scriptsize

\noindent\textbf{Instructions: }\vspace{-1mm}\begin{verbatim}

1. It is prohibited to search books and notes.  2. The use of electronic devices is prohibited. 3. It will be
considered only the answers in "Part 1" area of the Multiple-Choice Questions page.
\end{verbatim}
\begin{center}\textbf{Parte 2 - Multiple-Choice Questions}\end{center}
{\small
\begin{questions}
\itemsep0pt\parskip0pt\parsep0pt
\question pergunta f�cil Q1b2 - com segunda varia��es da subclasse b:
\begin{choices}
\itemsep0pt\parskip0pt\parsep0pt
\choice resposta 1b2-a
\choice resposta 1b2-e
\choice resposta 1b2-d
\choice resposta 1b2-b
\choice resposta 1b2-c
\end{choices}
\question pergunta Q14
\begin{choices}
\itemsep0pt\parskip0pt\parsep0pt
\choice resposta 14e
\choice resposta 14b
\choice resposta 14a
\choice resposta 14d
\choice resposta 14c
\end{choices}
\question pergunta Q2  com exemplo de f�rmula em tex: 
 $\sin A \cos B = \frac{1}{2}\left[ \sin(A-B)+\sin(A+B) \right]$
\begin{choices}
\itemsep0pt\parskip0pt\parsep0pt
\choice resposta 2b
\choice resposta 2a
\choice resposta 2e
\choice resposta 2d
\choice resposta 2c
\end{choices}
\question pergunta Q13
\begin{choices}
\itemsep0pt\parskip0pt\parsep0pt
\choice resposta 13b
\choice resposta 13d
\choice resposta 13e
\choice resposta 13c
\choice resposta 13a
\end{choices}
\question pergunta Q6
\begin{choices}
\itemsep0pt\parskip0pt\parsep0pt
\choice resposta 6e
\choice resposta 6c
\choice resposta 6d
\choice resposta 6b
\choice resposta 6a
\end{choices}
\question pergunta Q5
\begin{choices}
\itemsep0pt\parskip0pt\parsep0pt
\choice resposta 5d
\choice resposta 5e
\choice resposta 5b
\choice resposta 5a
\choice resposta 5c
\end{choices}
\question pergunta Q12
\begin{choices}
\itemsep0pt\parskip0pt\parsep0pt
\choice resposta 12d
\choice resposta 12e
\choice resposta 12c
\choice resposta 12b
\choice resposta 12a
\end{choices}
\question pergunta Q11
\begin{choices}
\itemsep0pt\parskip0pt\parsep0pt
\choice resposta 11d
\choice resposta 11b
\choice resposta 11a
\choice resposta 11e
\choice resposta 11c
\end{choices}
\question pergunta Q10
\begin{choices}
\itemsep0pt\parskip0pt\parsep0pt
\choice resposta 0b
\choice resposta 0e
\choice resposta 0d
\choice resposta 0c
\choice resposta 0a
\end{choices}
\end{questions}
}
 \ \ \ 
 \newpage
\begin{center}\textbf{Parte 3 - Textual Questions
}\end{center}
{\small
\makeatletter\renewcommand*\cleardoublepage{\ifodd\c@page \else\hbox{}\newpage\fi}
\makeatother
\cleardoublepage
\begin{table}[h]\centering
\begin{tabular}{|p{16mm}|p{16cm}|}
\hline\multirow{4}{*}{\hspace{-2mm}\includegraphics[width=2cm]{../../../figs/ufabc.eps}}
&\vspace{-2mm}\noindent\large\textbf{Universidade Federal do ABC}\\
&\noindent\textbf{Course:} Processamento da Informa��o - BC0505\\
&\noindent\textbf{Teacher(s):} Francisco de Assis Zampirolli\\
&\noindent\textbf{Period:} 2/2016\hfill\textbf{Modality:} Presencial\hfill\textbf{Date:} 22/06/2016\hspace{-8mm}\\
\hline\end{tabular}
\end{table}
\vspace{-4mm}\small{
\noindent\textbf{Student:}  Fulano Junior
\hfill\textbf{Registration:} 11000123
\hfill\textbf{Room:} 2015$\_$BC0505$\_$q3$\_$A2
\hspace{-1mm}

\vspace{8mm}
\noindent\textbf{Sig.:}\rule{11.5cm}{0.1pt}\hfill\hspace{1cm}
}

\vspace{4mm}\noindent Escreva 222 um programa para inverter os elementos �mpares que est�o nas posi��es pares de um vetor com X elementos, onde X � um inteiro definido pelo usu�rio.
 
 \drawLines{10} 


 \ \ \ 
 \newpage
\makeatletter\renewcommand*\cleardoublepage{\ifodd\c@page \else\hbox{}\newpage\fi}
\makeatother
\cleardoublepage
\begin{table}[h]\centering
\begin{tabular}{|p{16mm}|p{16cm}|}
\hline\multirow{4}{*}{\hspace{-2mm}\includegraphics[width=2cm]{../../../figs/ufabc.eps}}
&\vspace{-2mm}\noindent\large\textbf{Universidade Federal do ABC}\\
&\noindent\textbf{Course:} Processamento da Informa��o - BC0505\\
&\noindent\textbf{Teacher(s):} Francisco de Assis Zampirolli\\
&\noindent\textbf{Period:} 2/2016\hfill\textbf{Modality:} Presencial\hfill\textbf{Date:} 22/06/2016\hspace{-8mm}\\
\hline\end{tabular}
\end{table}
\vspace{-4mm}\small{
\noindent\textbf{Student:}  Fulano Junior
\hfill\textbf{Registration:} 11000123
\hfill\textbf{Room:} 2015$\_$BC0505$\_$q3$\_$A2
\hspace{-1mm}

\vspace{8mm}
\noindent\textbf{Sig.:}\rule{11.5cm}{0.1pt}\hfill\hspace{1cm}
}

\vspace{4mm}\noindent O programa abaixo l� dois valores para as vari�veis X e Y, efetua a troca dos valores de forma que a vari�vel X passe a ter o valor de Y, e que a vari�vel Y passe a ter o valor de X. Complete a(s) instru��o(�es) "AQUI".
 \begin{verbatim}
 programa
 {
     funcao inicio()
     {
           real X, Y, aux
           leia (X, Y)
           AQUI
           escreva(X, Y)
      }
 }
 \end{verbatim}
 
 \drawLines{10} 


 \ \ \ 
 \newpage
}
\makeatletter\renewcommand*\cleardoublepage{\ifodd\c@page \else\hbox{}\newpage\fi}
\makeatother
\cleardoublepage
\begin{table}[h]\centering
\begin{tabular}{|p{16mm}|p{16cm}|}
\hline\multirow{4}{*}{\hspace{-2mm}\includegraphics[width=2cm]{../../../figs/ufabc.eps}}
&\vspace{-2mm}\noindent\large\textbf{Universidade Federal do ABC}\\
&\noindent\textbf{Course:} Processamento da Informa��o - BC0505\\
&\noindent\textbf{Teacher(s):} Francisco de Assis Zampirolli\\
&\noindent\textbf{Period:} 2/2016\hfill\textbf{Modality:} Presencial\hfill\textbf{Date:} 22/06/2016\hspace{-8mm}\\
\hline\end{tabular}
\end{table}
\vspace{-4mm}\small{
\noindent\textbf{Student:}  Gustavo Neto
\hfill\textbf{Registration:} 11000111
\hfill\textbf{Room:} 2015$\_$BC0505$\_$q3$\_$A2
\hspace{-1mm}

\vspace{8mm}
\noindent\textbf{Sig.:}\rule{11.5cm}{0.1pt}\hfill\hspace{1cm}
}\begin{pspicture}(6,0in)
\psbarcode[scalex=1.6,scaley=0.35]{000011000111}{}{ean13}
\end{pspicture}
\\{\scriptsize

\noindent\textbf{Instructions: }\vspace{-1mm}\begin{verbatim}

1. Paint only INSIDE OF THE CIRCLES of each question.  2. No Rasure. 3. Each question has only one correct answer.
4. It will be considered only the answers in "Part 1" area on this page to the questions of multiple choice.
\end{verbatim}
\begin{center}\textbf{Parte 1 - Answer Sheet - Do not use this sheet as a draft!}\end{center}
\vspace{-5mm}\noindent\varhrulefill[0.4mm]\vspace{1mm}

\vspace{-3mm}\noindent\varhrulefill[0.4mm]\vspace{1mm}

\begin{center}
\begin{tikzpicture}[font=\tiny]
  \foreach \letter/\position in {A/1,B/2,C/3,D/4,E/5} {
    \node[inner sep=3pt] at ({\position * 0.5},0) {\letter};
  }
  \foreach \line in {1,...,9} {
     \begin{scope}[xshift=0cm,yshift=-(\line-1+1)*5mm]
       \foreach \letter/\position in {A/1,B/2,C/3,D/4,E/5} {
           \node at (-0.1,0) {\line};
           \node[fill=black!30,draw,circle,inner sep=3pt] at ({\position * 0.5},0) {};
           \node[fill=white,draw,circle,inner sep=2pt] at ({\position * 0.5},0) {};
       }
     \end{scope}
  }
\end{tikzpicture}\hspace{1 cm}


\end{center}
\vspace{1cm}\noindent\varhrulefill[0.4mm]\vspace{1mm}

\vspace{-3mm}\noindent\varhrulefill[0.4mm]\vspace{1mm}


\begin{table}[h]
\centering
\textbf{Calculation of concepts} \\ \vspace{5mm}
\begin{tabular}{|c|c|c|c|} \hline
Questions 1-12 & Textual Questions & Final\\ \hline
& &  \\
& &  \\
& &  \\ \hline
\end{tabular}
\end{table}\newpage\thispagestyle{empty}\mbox{}
 \ \ 
\newpage
\makeatletter\renewcommand*\cleardoublepage{\ifodd\c@page \else\hbox{}\newpage\fi}
\makeatother
\cleardoublepage
\begin{table}[h]\centering
\begin{tabular}{|p{16mm}|p{16cm}|}
\hline\multirow{4}{*}{\hspace{-2mm}\includegraphics[width=2cm]{../../../figs/ufabc.eps}}
&\vspace{-2mm}\noindent\large\textbf{Universidade Federal do ABC}\\
&\noindent\textbf{Course:} Processamento da Informa��o - BC0505\\
&\noindent\textbf{Teacher(s):} Francisco de Assis Zampirolli\\
&\noindent\textbf{Period:} 2/2016\hfill\textbf{Modality:} Presencial\hfill\textbf{Date:} 22/06/2016\hspace{-8mm}\\
\hline\end{tabular}
\end{table}
\vspace{-4mm}\small{
\noindent\textbf{Student:}  Gustavo Neto
\hfill\textbf{Registration:} 11000111
\hfill\textbf{Room:} 2015$\_$BC0505$\_$q3$\_$A2
\hspace{-1mm}

\vspace{8mm}
\noindent\textbf{Sig.:}\rule{11.5cm}{0.1pt}\hfill\hspace{1cm}
}
\vspace{4mm}
\\{\scriptsize

\noindent\textbf{Instructions: }\vspace{-1mm}\begin{verbatim}

1. It is prohibited to search books and notes.  2. The use of electronic devices is prohibited. 3. It will be
considered only the answers in "Part 1" area of the Multiple-Choice Questions page.
\end{verbatim}
\begin{center}\textbf{Parte 2 - Multiple-Choice Questions}\end{center}
{\small
\begin{questions}
\itemsep0pt\parskip0pt\parsep0pt
\question pergunta Q3
\begin{choices}
\itemsep0pt\parskip0pt\parsep0pt
\choice resposta 3d
\choice resposta 3e
\choice resposta 3b
\choice resposta 3c
\choice resposta 3a
\end{choices}
\question pergunta Q2  com exemplo de f�rmula em tex: 
 $\sin A \cos B = \frac{1}{2}\left[ \sin(A-B)+\sin(A+B) \right]$
\begin{choices}
\itemsep0pt\parskip0pt\parsep0pt
\choice resposta 2b
\choice resposta 2d
\choice resposta 2a
\choice resposta 2e
\choice resposta 2c
\end{choices}
\question pergunta Q14
\begin{choices}
\itemsep0pt\parskip0pt\parsep0pt
\choice resposta 14e
\choice resposta 14a
\choice resposta 14d
\choice resposta 14c
\choice resposta 14b
\end{choices}
\question pergunta f�cil Q1a2 - com segunda varia��es da subclasse a:
\begin{choices}
\itemsep0pt\parskip0pt\parsep0pt
\choice resposta 1a2-b
\choice resposta 
 1a2-a
\choice resposta 1a2-c
\choice resposta 1a2-d
\choice resposta 1a2-e
\end{choices}
\question pergunta Q4
\begin{choices}
\itemsep0pt\parskip0pt\parsep0pt
\choice resposta 4a
\choice resposta 4c
\choice resposta 4e
\choice resposta 4d
\choice resposta 4b
\end{choices}
\question pergunta Q6
\begin{choices}
\itemsep0pt\parskip0pt\parsep0pt
\choice resposta 6c
\choice resposta 6b
\choice resposta 6e
\choice resposta 6a
\choice resposta 6d
\end{choices}
\question pergunta Q5
\begin{choices}
\itemsep0pt\parskip0pt\parsep0pt
\choice resposta 5b
\choice resposta 5a
\choice resposta 5e
\choice resposta 5d
\choice resposta 5c
\end{choices}
\question pergunta Q7
\begin{choices}
\itemsep0pt\parskip0pt\parsep0pt
\choice resposta 7e
\choice resposta 7d
\choice resposta 7a
\choice resposta 7b
\choice resposta 7c
\end{choices}
\question pergunta Q9
\begin{choices}
\itemsep0pt\parskip0pt\parsep0pt
\choice resposta 9d
\choice resposta 9c
\choice resposta 9a
\choice resposta 9e
\choice resposta 9b
\end{choices}
\end{questions}
}
 \ \ \ 
 \newpage
\begin{center}\textbf{Parte 3 - Textual Questions
}\end{center}
{\small
\makeatletter\renewcommand*\cleardoublepage{\ifodd\c@page \else\hbox{}\newpage\fi}
\makeatother
\cleardoublepage
\begin{table}[h]\centering
\begin{tabular}{|p{16mm}|p{16cm}|}
\hline\multirow{4}{*}{\hspace{-2mm}\includegraphics[width=2cm]{../../../figs/ufabc.eps}}
&\vspace{-2mm}\noindent\large\textbf{Universidade Federal do ABC}\\
&\noindent\textbf{Course:} Processamento da Informa��o - BC0505\\
&\noindent\textbf{Teacher(s):} Francisco de Assis Zampirolli\\
&\noindent\textbf{Period:} 2/2016\hfill\textbf{Modality:} Presencial\hfill\textbf{Date:} 22/06/2016\hspace{-8mm}\\
\hline\end{tabular}
\end{table}
\vspace{-4mm}\small{
\noindent\textbf{Student:}  Gustavo Neto
\hfill\textbf{Registration:} 11000111
\hfill\textbf{Room:} 2015$\_$BC0505$\_$q3$\_$A2
\hspace{-1mm}

\vspace{8mm}
\noindent\textbf{Sig.:}\rule{11.5cm}{0.1pt}\hfill\hspace{1cm}
}

\vspace{4mm}\noindent Escreva 111 um programa para inverter os elementos pares que est�o nas posi��es �mpares de um vetor com X elementos, onde X � um inteiro definido pelo usu�rio.
 
 \drawLines{10} 


 \ \ \ 
 \newpage
\makeatletter\renewcommand*\cleardoublepage{\ifodd\c@page \else\hbox{}\newpage\fi}
\makeatother
\cleardoublepage
\begin{table}[h]\centering
\begin{tabular}{|p{16mm}|p{16cm}|}
\hline\multirow{4}{*}{\hspace{-2mm}\includegraphics[width=2cm]{../../../figs/ufabc.eps}}
&\vspace{-2mm}\noindent\large\textbf{Universidade Federal do ABC}\\
&\noindent\textbf{Course:} Processamento da Informa��o - BC0505\\
&\noindent\textbf{Teacher(s):} Francisco de Assis Zampirolli\\
&\noindent\textbf{Period:} 2/2016\hfill\textbf{Modality:} Presencial\hfill\textbf{Date:} 22/06/2016\hspace{-8mm}\\
\hline\end{tabular}
\end{table}
\vspace{-4mm}\small{
\noindent\textbf{Student:}  Gustavo Neto
\hfill\textbf{Registration:} 11000111
\hfill\textbf{Room:} 2015$\_$BC0505$\_$q3$\_$A2
\hspace{-1mm}

\vspace{8mm}
\noindent\textbf{Sig.:}\rule{11.5cm}{0.1pt}\hfill\hspace{1cm}
}

\vspace{4mm}\noindent O programa abaixo l� dois valores para as vari�veis X e Y, efetua a troca dos valores de forma que a vari�vel X passe a ter o valor de Y, e que a vari�vel Y passe a ter o valor de X. Complete a(s) instru��o(�es) "AQUI".
 \begin{verbatim}
 programa
 {
     funcao inicio()
     {
           real X, Y, aux
           leia (X, Y)
           AQUI
           escreva(X, Y)
      }
 }
 \end{verbatim}
 
 \drawLines{10} 


 \ \ \ 
 \newpage
}
\end{document}